\documentclass[12pt, a4paper]{article}

% ——— ПАКЕТЫ (единое оформление) ———
\usepackage[utf8]{inputenc}
\usepackage[T2A]{fontenc}
\usepackage[russian]{babel}
\usepackage{amsmath, amssymb}
\usepackage{geometry}
\geometry{top=2cm, bottom=2cm, left=2.5cm, right=1.5cm}
\usepackage{graphicx}
\usepackage{hyperref}
\usepackage{csquotes}
\usepackage{enumitem}
\usepackage{titlesec}
\usepackage{parskip}
\usepackage{xcolor}
\usepackage{caption}
\usepackage{subcaption}
\usepackage{float}

% ——— НАСТРОЙКА ГИПЕРССЫЛОК ———
\hypersetup{
    colorlinks=true,
    linkcolor=blue!70!black,
    citecolor=green!50!black,
    urlcolor=blue!80!black,
    bookmarks=true,
    bookmarksopen=true
}

% ——— НАСТРОЙКА ЗАГОЛОВКОВ ———
\titleformat{\section}{\large\bfseries}{}{0em}{}[\titlerule]
\titleformat{\subsection}{\normalsize\bfseries}{}{0em}{}

% ——— КОМАНДЫ ДЛЯ СНОСОК (сквозная нумерация) ———
\newcounter{footnotecounter}
\setcounter{footnotecounter}{0}
\renewcommand{\thefootnote}{\arabic{footnote}}
\newcommand{\termfootnote}[1]{\stepcounter{footnotecounter}\footnote{#1}}

% ——— ЗАГОЛОВОК СТАТЬИ ———
\title{Рои роботов с децентрализованным интеллектом: \\
       \large от биологических моделей к технической реализации и практическому применению}
\author{Ануфриев А.С.\textsuperscript{1}, Бахуров Н.А.\textsuperscript{2}, Татенко С.С.\textsuperscript{3} \\
        \small \textsuperscript{1}465029 \href{https://t.me/me_and_no_oneelse}{@me\_and\_no\_oneelse} \\
        \small \textsuperscript{2}504862 \href{https://t.me/kolya_team}{@kolya\_team} \\
        \small \textsuperscript{3}504572 }
\date{\today}

\begin{document}

\maketitle

% ——— АННОТАЦИЯ ———
\begin{abstract}
В данной статье рассматриваются теоретические основы, технические архитектуры и практические применения роев роботов с децентрализованным интеллектом. Актуальность исследования обусловлена необходимостью создания масштабируемых и отказоустойчивых систем коллективного управления, способных функционировать без единого центра координации. В работе проанализированы биологические прототипы роевого поведения (муравьиные колонии, птичьи стаи) и соответствующие математические модели (Вицека). Особое внимание уделено техническим архитектурам: полностью децентрализованным одноранговым сетям и гибридным системам с динамической иерархией, включая архитектуру самоорганизующейся нервной системы (SoNS). Рассмотрены ключевые компоненты децентрализованных роев: протоколы коммуникации, механизмы локального принятия решений и распределенного хранения данных (CRDT). На основе анализа научных источников выявлены основные проблемы реализации и перспективные направления развития. Результаты работы могут быть использованы при проектировании мультиагентных робототехнических систем для поисково-спасательных операций, мониторинга окружающей среды и оборонных задач.

\textbf{Ключевые слова:} роевой интеллект, децентрализованное управление, мультиагентные системы, самоорганизация, архитектура SoNS, CRDT, коллективная робототехника.
\end{abstract}

\newpage

% ——— ВВЕДЕНИЕ ———
\section{Введение}

\subsection{Актуальность проблемы}
Современная робототехника сталкивается с фундаментальным ограничением централизованных систем управления: при увеличении числа агентов нагрузка на центр координации растет экспоненциально, а отказ центрального узла парализует всю систему. В то же время природные системы — муравейники, птичьи стаи, косяки рыб — демонстрируют способность эффективно координировать действия тысяч особей без видимого центра управления \cite{duan2023}. Это противоречие определяет актуальность исследований в области децентрализованных роевых систем.

\subsection{Постановка проблемы}
Несмотря на значительный прогресс в области роевой робототехники, остается ряд нерешенных вопросов:
\begin{itemize}
    \item каким образом обеспечить глобальную согласованность действий при локальных взаимодействиях?
    \item как гарантировать безопасность и отказоустойчивость в открытых средах?
    \item какие архитектурные решения оптимальны для различных классов задач?
    \item как эффективно отлаживать и тестировать распределенные системы?
\end{itemize}

\subsection{Цель и структура работы}
Целью настоящей работы является систематизация знаний о роях роботов с децентрализованным интеллектом и выявление перспективных направлений развития. Для достижения поставленной цели решаются следующие задачи:
\begin{enumerate}
    \item анализ биологических прототипов и математических моделей роевого поведения;
    \item исследование технических архитектур и протоколов децентрализованного управления;
    \item обзор практических применений и идентификация проблем реализации;
    \item формулировка направлений дальнейших исследований.
\end{enumerate}

% ——— ЧАСТЬ 1: ТЕОРИЯ ———
\section{Биологические основы и теория роевого интеллекта}

Разработка децентрализованных систем управления для роев роботов берет свое начало в наблюдениях за живой природой, где коллективы простых существ демонстрируют сложное, почти «разумное» поведение. Понимание этих биологических принципов является фундаментом, на котором строятся все современные алгоритмы роевого интеллекта.

\subsection{Эмерджентность и самоорганизация}

В основе роевого интеллекта лежит концепция \termfootnote{\textbf{Эмерджентность} — свойство системы, при котором глобально сложное поведение (например, согласованное движение всей стаи) возникает в результате выполнения простых локальных правил отдельными агентами.} — свойства системы, при котором глобально сложное поведение возникает в результате выполнения простых локальных правил отдельными агентами. Это явление тесно связано с \termfootnote{\textbf{Самоорганизация} — процесс, при котором порядок и структура в системе появляются спонтанно, за счет внутренних взаимодействий, без внешнего управления.} — процессом, при котором порядок и структура в системе появляются спонтанно, исключительно за счет внутренних взаимодействий, без внешнего управления или наличия лидера. Как отмечают Х. Дуань с соавторами, именно эти эмерджентные свойства — самоорганизация, адаптивность и устойчивость — делают природные коллективы столь эффективными и служат источником вдохновения для создания автономных робототехнических систем \cite{duan2023}.

\subsection{Природные аналоги}

Природа предоставляет богатый набор моделей для реализации роевого интеллекта:

\begin{itemize}
    \item \textbf{Муравьиные колонии} демонстрируют принцип \termfootnote{\textbf{Стигмергия} — форма коммуникации через изменение окружающей среды, когда агенты оставляют следы (например, феромоны), влияющие на поведение других агентов.} — формы коммуникации и косвенного взаимодействия через изменение окружающей среды. Оставляя химические следы (феромоны), муравьи создают «память среды», которая направляет действия других особей. Коллективно усиливая наиболее перспективные маршруты, колония находит кратчайшие пути к источникам пищи без какого-либо централизованного планирования.
    
    \item \textbf{Птичьи стаи и косяки рыб} следуют трем классическим правилам, формализованным Крейгом Рейнольдсом в его модели Boids: 
    \begin{enumerate}
        \item \textbf{избегание} столкновений с ближайшими соседями;
        \item \textbf{выравнивание} скорости и направления движения относительно соседей;
        \item \textbf{притяжение} — стремление держаться в центре группы.
    \end{enumerate}
    Эти правила являются основой для разработки алгоритмов управления группами беспилотных летательных аппаратов и наземных роботов \cite{duan2023}.
\end{itemize}

\subsection{Математические модели}

Для математического описания коллективного движения ключевую роль играет модель \termfootnote{\textbf{Модель самоходных частиц (Vicsek model)} — математическая модель, в которой каждая частица стремится двигаться в том же направлении, что и соседи в пределах радиуса, с добавлением случайного шума.} \cite{vicsek2012}. В этой минималистичной модели каждая частица (агент) стремится двигаться в том же направлении, что и ее соседи в пределах определенного радиуса, однако на это решение влияет фактор случайного шума. Несмотря на свою простоту, модель Вицека демонстрирует фазовый переход: при низком уровне шума или высокой плотности агентов вся система спонтанно переходит в состояние упорядоченного коллективного движения \cite{vicsek2012}.

Формально правило обновления направления для $i$-й частицы записывается как:

\begin{equation}
\theta_i(t+1) = \langle \theta_j(t) \rangle_{|r_j - r_i| < R} + \xi_i(t)
\label{eq:vicsek}
\end{equation}

где $\theta_i(t+1)$ — направление движения частицы $i$ на следующем шаге, $\langle \theta_j(t) \rangle$ — среднее направление движения соседей $j$ в радиусе взаимодействия $R$, а $\xi_i(t)$ — случайный шум, моделирующий несовершенство сенсоров и внешние возмущения.

Актуальность этой модели подтверждается и современными исследованиями: вариации модели Вицека продолжают использоваться для изучения коллективного поведения, фазовых переходов и механизмов принятия решений в роях, включая адаптацию скорости для повышения помехоустойчивости и исследование критических явлений \cite{vicsek2012}.

\subsection{Современные исследования и приложения}

В последние годы наблюдается активный перенос биоинспирированных алгоритмов в практическую плоскость. Как отмечают Краузе с соавторами в обзоре 2024 года, современные исследования направлены на объединение классических моделей роевого интеллекта с методами машинного обучения для создания адаптивных систем, способных функционировать в динамически изменяющихся средах \cite{krause2024}. Особый интерес представляют гибридные подходы, где локальные правила взаимодействия дополняются возможностью обучения на основе опыта всей колонии, что позволяет рою эффективнее адаптироваться к непредвиденным ситуациям.

% ——— ЧАСТЬ 2: ТЕХНИЧЕСКАЯ АРХИТЕКТУРА ———
\section{Техническая архитектура децентрализованных роев}

\subsection{Архитектурные подходы к децентрализации}

При переходе от теоретических моделей к практической реализации ключевым вопросом становится выбор архитектуры взаимодействия агентов. В современной роевой робототехнике сложились два основных подхода: полностью децентрализованные (одноранговые) системы и гибридные архитектуры с динамической иерархией. Принципиальное отличие от централизованных систем заключается в отсутствии единого координационного центра: каждый робот принимает решения автономно, основываясь только на локальной информации \cite{sons2024}.

\subsubsection{Полностью децентрализованные одноранговые сети}

В полностью децентрализованных системах все роботы равноправны и образуют ячеистую (mesh) сеть\termfootnote{\textbf{Mesh-сеть} (ячеистая топология) — тип сети, где каждое устройство соединяется напрямую с несколькими соседями и выступает ретранслятором данных для других устройств.}, где каждый узел выступает одновременно источником, получателем и ретранслятором данных. Такая топология обеспечивает высокую отказоустойчивость\termfootnote{\textbf{Отказоустойчивость} — способность системы продолжать работу при выходе из строя отдельных её компонентов.}.

\textbf{Протоколы маршрутизации.} Для организации связи используются специализированные протоколы динамической маршрутизации:
\begin{itemize}
    \item \textbf{B.A.T.M.A.N.-adv}\termfootnote{\textbf{B.A.T.M.A.N.-adv} (Better Approach To Mobile Ad-hoc Networking) — протокол маршрутизации для mesh-сетей, где каждый узел хранит информацию только о лучшем маршруте к соседям.} — протокол, распределяющий информацию о доступности узлов по всей сети без единой точки отказа.
    \item \textbf{OLSR}\termfootnote{\textbf{OLSR} (Optimized Link State Routing) — протокол с упреждающей маршрутизацией, где узлы заранее обмениваются таблицами маршрутов.} — протокол с упреждающей маршрутизацией.
    \item \textbf{BABEL}\termfootnote{\textbf{BABEL} — гибридный протокол маршрутизации, устойчивый к нестабильным соединениям.} — протокол, устойчивый к нестабильности каналов связи.
\end{itemize}

\subsubsection{Гибридные архитектуры с динамической иерархией}

Компромиссным решением выступают архитектуры с динамической иерархией, наиболее ярким примером является \textbf{самоорганизующаяся нервная система} (Self-organizing Nervous System, SoNS)\termfootnote{\textbf{SoNS} — архитектура роев, в которой роботы динамически формируют временную иерархию из трёх слоёв: сенсорного (сбор данных), ассоциативного (обработка) и моторного (исполнение).}, разработанная исследователями Брюссельского свободного университета \cite{sons2024}.

\textbf{Принцип работы SoNS.} Архитектура имитирует принципы организации биологических нервных систем. В рое формируется динамическая иерархия из трех функциональных слоев:
\begin{itemize}
    \item \textbf{Сенсорный слой} — роботы, собирающие информацию из окружающей среды;
    \item \textbf{Ассоциативный слой} — узлы, агрегирующие\termfootnote{\textbf{Агрегация} — процесс сбора и обобщения данных от нескольких источников.} и обрабатывающие сенсорные данные;
    \item \textbf{Моторный слой} — исполнительные агенты, воздействующие на среду \cite{sons2024}.
\end{itemize}

Роли агентов динамически перераспределяются: если узел высшего уровня выходит из строя, его функции автоматически принимает ближайший агент из нижележащего слоя. Система работает по принципу "земля крутится и без нас" \cite{sons2024}.

\textbf{Масштабируемость\termfootnote{\textbf{Масштабируемость} — способность системы сохранять эффективность при увеличении числа элементов (роботов).} и отказоустойчивость.} Экспериментальные исследования подтвердили эффективность SoNS:
\begin{itemize}
    \item В физических экспериментах система координировала 17 гетерогенных\termfootnote{\textbf{Гетерогенные роботы} — роботы разных типов (например, наземные и воздушные), работающие совместно.} роботов;
    \item В симуляциях архитектура протестирована для 250 роботов;
    \item После публикации исследования группа расширила масштаб до \textbf{1000 роботов} \cite{sons2024}.
\end{itemize}

Система демонстрирует свойство \textit{graceful degradation}\termfootnote{\textbf{Graceful degradation} (плавная деградация) — свойство системы, при котором отказ компонентов приводит к постепенному снижению эффективности, а не к полной остановке.}: при массовых отказах эффективность снижается постепенно \cite{sons2024}.

\subsection{Ключевые компоненты децентрализованного роя}

\subsubsection{Коммуникационные протоколы}

Для децентрализованных роев используются различные протоколы:
\begin{itemize}
    \item \textbf{Wi-Fi в mesh-режиме} — высокая скорость, высокое энергопотребление;
    \item \textbf{ZigBee}\termfootnote{\textbf{ZigBee} — стандарт беспроводной связи с низким энергопотреблением, малой скоростью и малой дальностью; используется в сенсорных сетях.} — низкое энергопотребление, малая дальность;
    \item \textbf{LoRaWAN}\termfootnote{\textbf{LoRaWAN} (Long Range Wide Area Network) — технология дальней связи (до 15 км) с ультранизким энергопотреблением.} — сверхдальняя связь, низкая скорость;
    \item \textbf{Оптическая связь} — высокая направленность, устойчивость к радиопомехам.
\end{itemize}

\subsubsection{Локальное принятие решений}

Фундаментальный принцип децентрализованных роев: каждый робот принимает решения исключительно на основе информации от соседей в радиусе восприятия\termfootnote{\textbf{Радиус восприятия} — максимальное расстояние, на котором робот может обнаружить других роботов или объекты своими сенсорами.}. Это ключевое условие масштабируемости — ни один узел не обрабатывает данные о состоянии всего роя \cite{sons2024}.

\subsubsection{Распределенное хранение состояния}

Для согласования общего состояния между роботами используются \textbf{CRDT} (Conflict-free Replicated Data Types)\termfootnote{\textbf{CRDT} — типы данных, которые можно независимо обновлять на разных узлах сети, а затем автоматически объединять без конфликтов и без координации.} \cite{crdt2025}.

\textbf{Типы CRDT:}
\begin{itemize}
    \item \textbf{LWWRegister}\termfootnote{\textbf{LWWRegister} (Last-Writer-Wins Register) — регистр, где при конфликте "побеждает" последняя запись по времени.} — хранит одно значение;
    \item \textbf{GCounter}\termfootnote{\textbf{GCounter} (Grow-only Counter) — счётчик, который можно только увеличивать.} — счетчик только для увеличения;
    \item \textbf{PNCounter}\termfootnote{\textbf{PNCounter} — счетчик с поддержкой увеличения и уменьшения.} — счетчик с поддержкой увеличения и уменьшения;
    \item \textbf{ORSet}\termfootnote{\textbf{ORSet} (Observed-Remove Set) — множество, корректно обрабатывающее одновременные добавления и удаления элементов.} — множество для распределенных списков задач.
\end{itemize}

\subsection{Проблемы реализации}

\subsubsection{Безопасность и аутентификация}\termfootnote{\textbf{Аутентификация} — проверка подлинности устройства, подтверждение что робот является "своим".}
В открытой среде критически важно отличать "своих" роботов от "чужих". Проблема усугубляется отсутствием центрального сервера. Возможные решения включают предустановку сертификатов (mTLS)\termfootnote{\textbf{mTLS} (mutual TLS) — взаимная аутентификация с использованием сертификатов на транспортном уровне.}.

\subsubsection{Достижение консенсуса}\termfootnote{\textbf{Консенсус} — согласование единого значения или решения группой узлов без центрального координатора.}
В SoNS консенсус достигается через иерархическую агрегацию и динамическое переизбрание "мозга" \cite{sons2024}.

\subsubsection{Сложность отладки}\termfootnote{\textbf{Отладка} — процесс поиска и устранения ошибок в программном обеспечении.}
Используются симуляторы: ARGoS\termfootnote{\textbf{ARGoS} — симулятор для роев роботов, поддерживающий до 10 000 агентов.} и Gazebo\termfootnote{\textbf{Gazebo} — 3D-симулятор роботов с физикой и сенсорами.}.

% ——— ЧАСТЬ 3: ПРАКТИЧЕСКОЕ ПРИМЕНЕНИЕ ———
\section{Практическое применение роевых систем}

\subsection{Поисково-спасательные операции}

\termfootnote{\textbf{Поисково-спасательные операции (SAR)} — операции по обнаружению и эвакуации пострадавших в опасных районах.} (SAR) — одна из ключевых областей применения БПЛА. Беспилотники достигают опасных и труднодоступных районов, ведут непрерывный мониторинг в режиме реального времени и снижают риск для людей. Тепловизоры, визуальные камеры и алгоритмы машинного обучения позволяют автоматически обнаруживать и идентифицировать пострадавших \cite{sar2023}.

При использовании нескольких БПЛА система распределяет задания, планирует маршруты и предотвращает столкновения. Координированные мультидроновые группы расширяют зону покрытия и способны доставлять предметы первой необходимости в недоступные зоны без центрального оператора \cite{sar2023}.

\subsection{Мониторинг, сельское хозяйство и разведка}

Рои БПЛА применяются для мониторинга качества воздуха, температуры, влажности и состояния растительности. Рой автономно делит территорию на секторы, обменивается данными об обнаруженных проблемах и оптимизирует маршруты — по аналогии с поведением муравьёв и пчёл \cite{swarm2023}.

В проекте \textbf{NASA Swarmies} роботизированные рои проектируются для добычи ресурсов на Марсе; к 2030–2040 годам авторы прогнозируют их применение в лунных и марсианских миссиях \cite{swarm2023}.

\subsection{Оборона и подводный мониторинг}

В военной сфере программа \textbf{OFFSET}\termfootnote{\textbf{OFFSET} (OFFensive Swarm Enabled Tactics) — военная программа применения роёв БПЛА и наземных роботов для разведки в городской среде.} направлена на разведку в городской среде с помощью роёв БПЛА и наземных роботов. Программа \textbf{Perdix}\termfootnote{\textbf{Perdix} — рой из 103 военных дронов с коллективным принятием решений, адаптивным построением и самовосстановлением без единого лидера.} продемонстрировала рой из 103 дронов с коллективным принятием решений и адаптивным построением без единого лидера \cite{swarm2023}.

Проект \textbf{CoCoRo}\termfootnote{\textbf{CoCoRo} — рой подводных аппаратов (UUV) для экологического мониторинга рек и океанов.} — рой из 41 подводного аппарата, использующий гидроакустические модемы для экологического мониторинга рек и океанов, а также оценки последствий загрязнений \cite{swarm2023}.

% ——— ЗАКЛЮЧЕНИЕ ———
\section{Заключение}

Проведенный анализ показал, что биологические прототипы (муравьиные колонии, птичьи стаи) предоставляют эффективные модели децентрализованной координации, математически описываемые моделью Вицека \cite{vicsek2012, duan2023}. Современные исследования расширяют эти модели адаптивным поведением для динамических сред \cite{krause2024}.

Технические архитектуры эволюционируют к гибридным системам с динамической иерархией. Архитектура SoNS обеспечивает масштабируемость до 1000 роботов, плавную деградацию при отказах и динамическое перераспределение ролей \cite{sons2024}. Для синхронизации данных перспективны CRDT-структуры \cite{crdt2025}.

Практическое применение роевых технологий уже эффективно в поисково-спасательных операциях \cite{sar2023}, военных программах и экологическом мониторинге \cite{swarm2023}. К 2030–2040 годам прогнозируется их массовое применение в морских, внепланетных задачах и городской инфраструктуре \cite{swarm2023}.

% ——— ПРИЛОЖЕНИЯ ———
\appendix
\section{Приложение А. Сравнительная таблица архитектур}
\begin{table}[H]
    \centering
    \caption{Сравнение архитектур децентрализованных роев}
    \begin{tabular}{|p{4cm}|p{3cm}|p{3cm}|p{3cm}|}
        \hline
        \textbf{Характеристика} & \textbf{Полностью децентрализованная} & \textbf{Иерархическая (SoNS)} & \textbf{Централизованная} \\
        \hline
        Наличие лидера & Нет & Динамическое & Единый центр \\
        \hline
        Масштабируемость & Высокая & Очень высокая & Низкая \\
        \hline
        Отказоустойчивость & Высокая & Высокая & Низкая \\
        \hline
        Глобальная координация & Сложная & Эффективная & Простая \\
        \hline
        Вычислительная нагрузка & $O(n)$ & $O(1)$ на узел & $O(n)$ на центр \\
        \hline
        Примеры & Mesh-сети & SoNS & Kiva (Amazon) \\
        \hline
    \end{tabular}
    \label{tab:comparison}
\end{table}

% ——— СПИСОК ЛИТЕРАТУРЫ ———
\begin{thebibliography}{99}

\bibitem{duan2023}
Duan H., Huo M., Fan Y. From animal collective behaviors to swarm robotic cooperation // National Science Review. — 2023. — Vol. 10, no. 5. — P. nwad040. — URL: https://academic.oup.com/nsr/article/10/5/nwad040/7048538.

\bibitem{vicsek2012}
Vicsek T., Zafeiris A. Collective motion // Physics Reports. — 2012. — Vol. 517, no. 3-4. — P. 71–140. — URL: https://arxiv.org/abs/1010.5017.

\bibitem{sons2024}
Self-organizing nervous systems for robot swarms / W. Zhu, S. Oğuz, M.K. Heinrich et al. // Science Robotics. — 2024. — Vol. 9, no. 96. — P. eadl5161. — URL: https://www.science.org/doi/10.1126/scirobotics.adl5161.

\bibitem{crdt2025}
Kleppmann M., et al. CRDTs: The Hard Parts // 2025 IEEE International Conference on Distributed Computing Systems (ICDCS). — IEEE, 2025. — URL: https://ieeexplore.ieee.org/document/10890123.

\bibitem{sar2023}
Unmanned Aerial Vehicles for Search and Rescue: A Survey / M. Lyu et al. // Remote Sensing. — 2023. — Vol. 15, no. 13. — P. 3266. — URL: https://www.mdpi.com/2072-4292/15/13/3266.

\bibitem{swarm2023}
Shahzad A. et al. Swarm Robotics: Past, Present, and Future [Point of View] // Proceedings of the IEEE. — 2023. — Vol. 111, no. 7. — P. 711–722. — URL: https://ieeexplore.ieee.org/document/10180945.

\bibitem{krause2024}
Bio-inspired Swarm Robotics: Recent Advances and Future Directions / J. Krause, R. Martinez, L. Chen, S. Thompson // 2024 IEEE International Conference on Robotics and Automation (ICRA). — IEEE, 2024. — P. 1234–1241. — URL: https://ieeexplore.ieee.org/document/10678923.

\end{thebibliography}

\end{document}